\scanpage{117}
Prenons $\mathscr X=\Delta_n$, alors \deleted{\unsure{les|le}} le diagramme (1) devient
\diagram{\begin{tikzcd}[ampersand replacement=\&,column sep=small]
 \mathbb R^*\Gamma(\Delta_n,S_k)\simeq S_k\arrow[r] \& \mathbb R^*\Gamma(\Delta_n,C^{**}\unsure{]|)}=\deleted{\unsure{A}}C_n^{**}\\
 \& C^{**}(\Delta_n)=C_n^{**}\arrow[u,"{\mathrm{id}_{C_n^{**}}}"']
\end{tikzcd}}
et la condition énoncée se réduit à celle concernant la flèche horizontale, qui signifie aussi
\diagram{\(\text{(3)}\quad\)\begin{tikzcd}[ampersand replacement=\&,column sep=0pt]
 H_C^{**}(\Delta_n)={} \& H^{**}(C_n^{**}) \& \deleted{\text{est isom}}\unsure{k|*}\\
 \& S_k\arrow[u,"{\varepsilon\ \text{augmentation}}"'] \&
\end{tikzcd}}
l'augmentation est un quasi-iso i.e.
\[
 \begin{aligned}
 &\deleted{\unsure{k}}\deleted{\unsure{=}}\deleted{\unsure{B}}S_k^i\longrightarrow
 \underset{\substack{\Vert\\\operatorname{Ker}(C_n^{i,0}\mathbin{-}C_n^{i1}).}}{H^{i0}(C_n^{**})}
 \deleted{\unsure{\simeq}}\ \text{est un}\\
 &\text{iso pour }i\text{ grand }(i\geq\lambda(\unsure{+0|+{\infty}|0|n})\ \text{disons})
 \end{aligned}\tag*{3 a)}
\]
\[
 \begin{aligned}
 &\forall p>0,\qquad H^{ip}(C_n^{**})=0\ \text{pour }i\\
 &i\text{ grand (dépendant de }p),\ \text{soit}\\
 &H^{ip}(\deleted{\unsure{S}}C_n^{**})=0\quad\text{si}\quad i\geq\deleted{\unsure{g}}\lambda(p)
 \end{aligned}\tag*{3 b)}
\]
où on peut supposer $\lambda:\mathbb N^{*}\mathbin{\unsure{-|\to}}\mathbb N^{*}$ strictement croissante.

A priori, ce $\lambda=\lambda_n$ dépend du choix de $n$. On suspecte qu'il faut pouvoir prendre $\lambda$ indépendant de $n$ --
% archive sheet number: 116
